\chapter{Lipoedema}
\index{Lipoedema}

\section*{Definition and Overview}
Lipoedema is a chronic, progressive disorder of adipose tissue characterised by the symmetrical, disproportionate accumulation of subcutaneous fat in the lower extremities, and sometimes the upper extremities, with strict sparing of the hands and feet. This condition almost exclusively affects females and is frequently triggered or exacerbated by periods of significant hormonal fluctuation, such as puberty, pregnancy, oral contraceptive use, or menopause. Lipoedema is characterised not only by physical disproportion but also by tissue tenderness, easy bruising, and resistance to standard weight loss measures such as diet and exercise. It is distinct from general obesity and primary lymphoedema, though advanced stages can lead to secondary lymphatic insufficiency, a condition known as lipo-lymphoedema.

\section*{Epidemiology}
Lipoedema is a widely underdiagnosed and misdiagnosed condition, often mistaken for simple lifestyle-induced obesity or bilateral lymphoedema. It is estimated to affect up to 10\% of the female population worldwide to some degree. The condition is exceptionally rare in males, in whom it is typically associated with severe hormonal imbalances, such as hypogonadism, Klinefelter syndrome, or advanced liver disease. Symptoms most commonly manifest during puberty or other hormonal transitions. In many countries, awareness of lipoedema has grown, but many patients still experience a diagnostic delay of several years or decades, which often results in progressive physical impairment and psychological distress.

\section*{Aetiology and Pathophysiology}
The exact pathophysiological mechanisms underlying lipoedema remain to be fully elucidated, but a combination of genetic, hormonal, and microvascular factors is involved.

\begin{itemize}
    \item \textbf{Genetic predisposition:} A strong familial occurrence is observed, often demonstrating an autosomal dominant inheritance pattern with variable penetrance.
    \item \textbf{Hormonal influence:} The female-specific nature of the disease indicates that oestrogen plays a key role in regulating adipocyte proliferation, distribution, and vascular permeability in the subcutis.
    \item \textbf{Microvascular and capillary fragility:} Chronic capillary leakiness leads to the accumulation of protein-rich fluid in the interstitial space. This chronic micro-oedema induces localized, low-grade inflammation, adipocyte hypertrophy, and progressive interstitial fibrosis.
    \item \textbf{Adipose tissue remodeling:} Subcutaneous fat cells undergo both hypertrophy (enlargement) and hyperplasia (proliferation), particularly in the thighs, calves, and buttocks, leading to a lobular, nodular tissue structure.
    \item \textbf{Lymphatic overload:} In early stages, lymphatic transport capacity is normal or even increased to compensate for microvascular leakage. Over time, the physical pressure of the hypertrophied fat tissue and chronic inflammation damages the lymphatic vessels, leading to secondary lymphatic insufficiency.
\end{itemize}

\section*{Clinical Features}
Lipoedema typically presents with a distinct pattern of physical and sensory symptoms that differentiate it from other causes of limb swelling.

\begin{itemize}
    \item Symmetrical, bilateral enlargement of the legs (from the hips to the ankles) or arms (from the shoulders to the wrists), with a sharp demarcation or "cuff" at the ankles or wrists, leaving the feet and hands unaffected.
    \item Pain, tenderness, and hyperalgesia (extreme sensitivity) in the subcutaneous fat, where even light pressure from clothing or touch can cause discomfort.
    \item Capillary fragility, resulting in easy, spontaneous bruising with minimal or no known trauma.
    \item A persistent sensation of heaviness, tightness, and aching in the affected limbs, which typically worsens at the end of the day, after prolonged standing, or in hot weather.
    \item Texture changes in the subcutaneous tissue, which initially feels soft and nodular (like "small beans" or "rice grains" under the skin) and progresses to large, lobular fat masses in advanced stages.
    \item Resistance to calorie-restricted diets and intensive exercise programmes, which fail to reduce the abnormal limb fat despite weight loss occurring in the torso.
\end{itemize}

\section*{Differential Diagnosis}
Distinguishing lipoedema from other conditions that cause limb enlargement is critical to ensuring appropriate management.

\begin{itemize}
    \item \textbf{Generalised obesity:} Fat accumulation occurs throughout the entire body, including the hands, feet, and face. It does not present with tissue tenderness, easy bruising, or the characteristic ankle/wrist cuff, and is responsive to lifestyle modifications.
    \item \textbf{Lymphoedema:} Typically asymmetrical and involves the dorsum of the feet or hands. It presents with pitting oedema in early stages, a positive Stemmer's sign (inability to pinch the skin on the second toe), and is not characteristically painful or prone to bruising.
    \item \textbf{Chronic venous insufficiency (CVI):} Characterised by venous hypertension, varicose veins, pitting oedema, skin hyperpigmentation (hemosiderin staining), and stasis dermatitis around the medial malleolus, though it can coexist with lipoedema.
    \item \textbf{Dercum's disease (adiposis dolorosa):} A rare condition characterized by painful lipomas or fatty deposits, but these are typically multiple, discrete, and distributed across the trunk and proximal limbs rather than presenting as symmetrical limb enlargement.
    \item \textbf{Lipodystrophy:} A group of disorders characterized by selective loss of adipose tissue, which is the opposite of the fat accumulation seen in lipoedema.
\end{itemize}

\section*{When to Seek Medical Care}
Patients should consult a general practitioner if they suspect they have lipoedema to establish a management plan and prevent progression.

\textbf{Contact your local emergency services for:}
\begin{itemize}
    \item Sudden, asymmetrical swelling, redness, warmth, or severe pain in one limb, which may indicate a deep vein thrombosis (DVT) or acute cellulitis.
    \item Chest pain, sudden shortness of breath, or coughing up blood, which may indicate a pulmonary embolism.
    \item Signs of systemic infection, such as high fever, chills, rapid heart rate, or confusion.
    \item Severe difficulty breathing or gasping for air.
\end{itemize}

\section*{Diagnosis and Investigations}
The diagnosis of lipoedema is primarily clinical, based on a detailed medical history and physical examination. There are currently no diagnostic blood tests or imaging studies specific to lipoedema.

\begin{itemize}
    \item \textbf{Physical examination:} Assessing the symmetry of fat accumulation, checking for sparing of the hands and feet, verifying a negative Stemmer's sign, and assessing tissue tenderness and capillary fragility.
    \item \textbf{Duplex ultrasound:} Performed to evaluate the venous system and rule out deep vein thrombosis or chronic venous insufficiency.
    \item \textbf{Lymphoscintigraphy:} Typically shows normal or slightly delayed lymphatic transport in early-to-mid stage lipoedema, helping to differentiate it from primary lymphoedema.
    \item \textbf{MRI or CT scan:} Can visualize the depth and symmetry of subcutaneous fat hypertrophy and rule out other structural pathology.
    \item \textbf{Thyroid and metabolic panels:} Performed to rule out endocrine or metabolic conditions that could contribute to weight gain and fluid retention.
\end{itemize}

\section*{Management}
Management is multi-disciplinary, focusing on symptom control, improving mobility, preventing complications, and supporting psychological health.

\begin{itemize}
    \item \textbf{Compression therapy:} Wearing flat-knit compression garments (which are stiffer and provide better containment than circular-knit garments) to support the tissues, reduce pain, and prevent secondary fluid accumulation.
    \item \textbf{Manual lymphatic drainage (MLD):} A gentle massage technique performed by a qualified lymphoedema therapist to stimulate lymphatic flow and alleviate aching and discomfort.
    \item \textbf{Anti-inflammatory diet:} Adopting eating patterns such as the Mediterranean diet or a low-carbohydrate/ketogenic diet to help reduce systemic inflammation and manage co-existing general obesity.
    \item \textbf{Low-impact physical activity:} Gentle exercise such as swimming, aqua aerobics, cycling, or walking, which utilises the muscle pump to promote lymphatic and venous return without placing excessive stress on joints.
    \item \textbf{Surgical intervention (Lymphatic-sparing liposuction):} Tumescent local anaesthesia (TLA) or water-jet assisted liposuction (WAL) can be performed by experienced surgeons to remove abnormal lipoedema tissue while protecting the lymphatic vessels.
    \item \textbf{Psychosocial support:} Counselling or cognitive behavioural therapy to address the significant body image distress, depression, and anxiety commonly associated with the condition.
\end{itemize}

\section*{Complications}
\begin{itemize}
    \item Lipo-lymphoedema, where chronic tissue pressure and microvascular damage lead to secondary failure of the lymphatic system, resulting in pitting fluid accumulation in the feet and hands.
    \item Recurrent skin infections, such as cellulitis or erysipelas, due to compromised local lymphatic clearance and skin fold irritation.
    \item Mobility impairment, joint degeneration, and gait abnormalities (particularly osteoarthritis of the knees and hips) due to the mechanical weight of the limbs.
    \item Severe psychological distress, including body dysmorphia, clinical depression, anxiety, social isolation, and eating disorders.
    \item Chronic pain and physical exhaustion, significantly reducing overall quality of life.
\end{itemize}

\section*{Prognosis and Outcomes}
Lipoedema is a progressive, lifelong condition that currently has no cure. Without appropriate management, the condition typically worsens over time, leading to significant mobility issues and secondary lymphatic complications. Conservative therapies (compression, manual lymphatic drainage, and exercise) are highly effective at controlling symptoms, reducing pain, and preventing the onset of lipo-lymphoedema, but they do not reverse the abnormal fat accumulation. Lymphatic-sparing liposuction offers the most durable outcome for reducing limb volume, relieving pain, and improving long-term mobility and quality of life.

\section*{Prevention and Self-Care}
\begin{itemize}
    \item Wear prescribed flat-knit compression garments consistently, especially during air travel, exercise, or prolonged standing.
    \item Maintain a regular schedule of gentle physical activity, prioritizing water-based exercises to utilize the natural compression of water.
    \item Maintain a healthy weight through anti-inflammatory dietary habits to avoid adding generalised obesity to lipoedema tissues.
    \item Practice meticulous skin care: wash and moisturise the limbs daily, and promptly treat any minor cuts or insect bites with antiseptic.
    \item Elevate the legs whenever resting to reduce gravitational swelling and ease the heavy sensation.
    \item Avoid restrictive footwear and clothing that could create tourniquet effects at the ankles, knees, or waist.
\end{itemize}

\section*{Sources and Further Reading}
The following reputable organizations provide further information, research, and support networks for individuals and healthcare professionals dealing with lipoedema.

\begin{itemize}
    \item Lipoedema Australia. \textit{Lipoedema diagnosis, resources, and patient support.} https://www.lipoedema.org.au/
    \item NHS. \textit{Lipoedema: symptoms, diagnosis, and treatment options.} https://www.nhs.uk/conditions/lipoedema/
    \item Lipoedema UK. \textit{Clinical consensus, research, and guidelines for managing lipoedema.} https://www.lipoedema.co.uk/
    \item Mayo Clinic. \textit{Lipoedema vs obesity: what is the difference?} https://www.mayoclinic.org/
    \item Healthdirect Australia. \textit{Lipoedema overview and conservative therapies.} https://www.healthdirect.gov.au/
\end{itemize}
